% =====================================================================
% Schematic of the parameters carried by a pulse descriptor word (PDW)
% for use in \cref{sec:bg:pdw}.  A short pulse train illustrates the
% time-domain PDW fields (TOA, PW, PRI, amplitude, RF carrier); a small
% inset shows AOA as a spatial bearing on the receive aperture.
% Loaded from sections/02_background/01_radar_systems_passive_reception.tex via \input{...}.
% =====================================================================
\begin{tikzpicture}[
    >=Latex,
    font=\small,
    pulse/.style={draw=blue!60!black, thick, fill=blue!10},
    rfwave/.style={very thin, red!60!black},
    annot/.style={font=\small\itshape},
    measure/.style={<->, thick, shorten >=1pt, shorten <=1pt},
    drop/.style={dashed, thin, gray!70!black},
  ]

  % --- pulse train geometry (TOA = leading edge, all PW equal here) ---
  \def\toaA{1.4}
  \def\toaB{4.4}
  \def\toaC{7.4}
  \def\pw{0.9}
  \def\ampA{2.4}
  \def\ampB{2.2}
  \def\ampC{2.5}

  % --- axes -----------------------------------------------------------
  \draw[->] (-0.2, 0) -- (10.4, 0) node[below] {time $t$};
  \draw[->] (0, -0.2) -- (0, 3.6) node[left]  {amplitude};

  % --- pulse envelopes ------------------------------------------------
  \draw[pulse] (\toaA, 0) rectangle (\toaA + \pw, \ampA);
  \draw[pulse] (\toaB, 0) rectangle (\toaB + \pw, \ampB);
  \draw[pulse] (\toaC, 0) rectangle (\toaC + \pw, \ampC);

  % --- illustrative RF carrier inside the middle pulse ---------------
  \draw[rfwave, samples=120, smooth, domain=\toaB:\toaB+\pw]
    plot (\x, {0.5*\ampB + 0.42*\ampB*sin(2880*(\x - \toaB))});

  % --- TOA drop lines and labels --------------------------------------
  \foreach \t/\lbl in {%
      \toaA/$t_n$,%
      \toaB/$t_{n+1}$,%
      \toaC/$t_{n+2}$%
    } {%
    \draw[drop] (\t, 0) -- (\t, -0.35);
    \node[annot, below] at (\t, -0.32) {\lbl};
  }
  \node[annot, below right=1pt and -2pt] at (\toaA, -0.32)
    {\scriptsize\hspace{0.15em}(TOA)};

  % --- amplitude arrow on the first pulse -----------------------------
  \draw[measure] (\toaA - 0.5, 0) -- (\toaA - 0.5, \ampA);
  \node[annot, left] at (\toaA - 0.5, \ampA/2) {$A_n$};

  % --- PW arrow above the third pulse ---------------------------------
  \draw[measure] (\toaC, \ampC + 0.3) -- (\toaC + \pw, \ampC + 0.3);
  \node[annot, above] at (\toaC + \pw/2, \ampC + 0.3) {PW};

  % --- PRI arrow between the first two leading edges ------------------
  \pgfmathsetmacro{\priY}{\ampA + 0.95}
  \draw[measure] (\toaA, \priY) -- (\toaB, \priY);
  \node[annot, above] at ({(\toaA + \toaB)/2}, \priY) {PRI $= 1/$PRF};

  % --- RF label with leader line into the middle-pulse oscillation ----
  \node[annot, text=red!60!black] (rflabel) at (6.4, 0.55*\ampB)
    {RF $f_n$};
  \draw[->, very thin, red!60!black]
    (rflabel.west) -- (\toaB + \pw - 0.05, 0.55*\ampB);

  % --- AOA inset (upper right): bearing on the receive aperture -------
  \begin{scope}[shift={(9.0, 2.55)}]
    % receive aperture
    \draw[thick] (-0.30, 0) -- (0.30, 0) -- (0, 0.45) -- cycle;
    \node[font=\footnotesize, below=1pt] at (0, 0) {Rx};
    % boresight reference
    \draw[drop] (0, 0.45) -- (0, 1.55);
    % incoming wavefront direction
    \draw[->, thick] (1.0, 1.40) -- (0.04, 0.50);
    % angle arc
    \draw[->, thin] (0, 1.15) arc (90:50:0.7);
    \node[annot] at (0.55, 1.00) {AOA $\theta_n$};
  \end{scope}

\end{tikzpicture}
